\section{Introduction}
\label{sec:introduction}

Low-light enhancement has made steady progress by improving RGB reconstruction quality. Retinex-inspired decompositions, zero-reference curve estimation, and transformer-based restoration models have all improved the visual and numerical fidelity of enhanced images~\cite{land1971retinex,wei2018deepretinex,zhang2019kind,guo2020zerodce,li2021zerodcepp,cai2023retinexformer}. Yet the experiments in this work begin from an inconvenient observation: image fidelity can improve while the model's local exposure field becomes physically less meaningful. On UEFB-v2, an A-gate variant reaches a higher PSNR than our centered E-shape model, but its gauge-free exposure-shape correlation is negative. A PSNR-only benchmark would therefore select a model whose exposure field points in the wrong spatial direction.

This observation matters because exposure-field models are usually motivated not only as black-box enhancers, but also as interpretable mechanisms for uneven illumination. If the field is not identifiable, then the method's physical story becomes fragile: a model may look better in RGB space while failing the very exposure interpretation it claims to provide. The central question of this paper is therefore not whether another network can beat the strongest restoration baseline. Instead, we ask: \emph{what does it mean for a local exposure field to be identifiable, and how should we evaluate such a field when RGB fidelity and field consistency disagree?}

We argue that the core difficulty is gauge ambiguity. A local exposure field is naturally defined in log-luminance space, but gradient-domain or Poisson-like constraints are insensitive to additive shifts. Thus, the absolute gauge of the field and its spatial shape are entangled. This motivates the decomposition $E=S+\mu$, where $S=E-\mathrm{mean}(E)$ is a gauge-free spatial shape and $\mu$ is the image-level gauge. Under this view, a model should not be judged only by absolute field error or output PSNR; it should be evaluated by how well it predicts shape, gauge, and RGB output separately.

To operationalize this idea, we propose GIR-Field, a gauge-identifiable exposure-field learning and evaluation framework. The method uses centered exposure-shape calibration to supervise spatial exposure structure while ignoring additive gauge shifts. The evaluation protocol, UEFB-G, reports output-only metrics for all methods and internal-field metrics only for methods with explicit exposure predictions. It further includes gauge perturbations that test whether a metric can distinguish global gauge shifts from local shape distortions.

Our full frozen-evidence audit supports the need for this reframing. Across 500 UEFB-v2 test images, 100 LOL-v2-real images, and 100 LOL-v2-synthetic images, centered E-shape calibration improves UEFB gauge-free shape correlation over joint training by +0.1581 with a 95\% bootstrap confidence interval of [0.1369, 0.1791] and FDR-corrected $q=2.18\times10^{-43}$. It also improves real PSNR by +0.6963 dB and synthetic PSNR by +0.6321 dB. Conversely, the PSNR-high A-gate model has negative exposure-shape correlation on all internal-field splits, showing that PSNR-only ranking can hide a physically wrong field.

We deliberately calibrate the claim boundary. GIR-Field does not outperform Retinexformer~\cite{cai2023retinexformer} in RGB fidelity, and we do not present it as a SOTA low-light restoration method. Retinexformer remains substantially stronger on LOL-v2-real and LOL-v2-synthetic. The contribution is instead methodological: we identify a gauge ambiguity in exposure-field learning, provide metrics and perturbations that separate gauge and shape errors, and show that a centered shape objective can produce statistically reliable field improvements without sacrificing paired fidelity.

Our contributions are:
\begin{itemize}
    \item We formalize local exposure-field learning as a gauge-identifiability problem, decomposing $E$ into a gauge-free spatial shape $S$ and image-level gauge $\mu$.
    \item We introduce centered exposure-shape calibration, a gauge-invariant objective that improves exposure-shape correlation and paired fidelity in a full local audit.
    \item We propose UEFB-G, a gauge-controlled evaluation protocol that separates output-only metrics from internal-field metrics and uses perturbations to distinguish global gauge shifts from local shape distortions.
    \item We provide a claim-calibrated failure audit: PSNR-only ranking can select a physically incoherent exposure field, and recoverability risk is predictable but not yet well calibrated.
\end{itemize}
